\section{Discussion}\label{sec:discussion}
\subsection{What the guarantee provides in practice}\label{sec:disc-buys}
SCC quantifies the reliability of a transferred RUL interval from physical parameters before any
target failure data exists. When the runtime diagnostic confirms similitude (\emph{holds}),
arbitrarily large differences in the dominant operating parameters, the $10.2\times$ life range in
the testbed, cost nothing in coverage, because the dimensionless scale absorbs them. When similitude
is incomplete, the certificate degrades by the physical departure $\lVert\Delta\boldsymbol\psi\rVert$
and the diagnostic flags the regime. This is a different bargain from weighted
conformal~\cite{tibshirani2019conformal}: SCC trades the need for target failure data, which a new
asset lacks, for a falsifiable physical assumption that the diagnostic makes checkable.

\subsection{Deployment and fault-tolerant use}\label{sec:disc-deploy}
The two by-products of the method are what make it usable on a real asset. The a-priori bound lets a
planner quote a worst-case coverage shortfall, and therefore a worst-case reliability of the
maintenance decision, before the target unit is commissioned. The runtime diagnostic acts as an
online monitor: when its premise is unsupported it returns \emph{indeterminate} or \emph{violated},
and the system falls back to conservative data-driven conformal rather than issuing an interval whose
confidence is silently wrong. The result is graceful degradation, the property a safety-relevant
prognostic function needs, instead of an unannounced loss of coverage.

\subsection{Limitations and validation roadmap}\label{sec:disc-limits}
Four limitations bound the claims. (i) The guarantee is conditional on the physical hypotheses
H1--H2 (\ref{app:proof}), which are falsifiable rather than proven. (ii) The certificate bounds the
gap for the modelled groups; unmodelled physics contributes a residual that the departure study
quantifies and the diagnostic detects, but that must be bounded a priori to use the certificate under
incomplete similitude. (iii) Validation is on a controlled simulation, and the single field benchmark
returns \emph{indeterminate} by design rather than a positive result. (iv) The bound is conservative.
The decisive next step is validation on adequately powered run-to-failure field data, where the
diagnostic returns \emph{holds} and the certificate can be checked against measured cross-regime
coverage; such data is the goal of ongoing work on an operating industrial process.
